\section{Loopless directed carrier, actions, and squarefree dictionaries}
\label{sec:model}

Let $V=[n]=\{0,\ldots,n-1\}$ and
\[
  \Omega_n=\{(i,j)\in V^2:i\neq j\}.
\]
A labelled loopless digraph is a subset $P\subseteq\Omega_n$, or equivalently
its indicator $x^P\in\{0,1\}^{\Omega_n}$.  The selected group $G\leq S_n$
acts by
\[
  g\cdot(i,j)=(g(i),g(j)),\qquad
  g\cdot P=\{g\cdot a:a\in P\}.
\]
The objects to be reconstructed are the orbits $G\backslash
\{0,1\}^{\Omega_n}$.  There are $218,368,9608,17824$ of them for the four
actions in Table~\ref{tab:exact-panel}.

For a nonempty support $T\subseteq\Omega_n$, set
\[
  m_T(x)=\prod_{a\in T}x_a,
  \qquad
  q^G_{[T]}(x)=\sum_{U\in G\cdot T}m_U(x),
\]
where the sum is over the distinct supports in the orbit.  On a Boolean
pattern $P$ this specializes to
\begin{equation}
  q^G_{[T]}(P)=\#\{U\in G\cdot T:U\subseteq P\}.
  \label{eq:contained-support-count}
\end{equation}
Thus every coordinate is an integer count.  It depends only on the orbit of
$P$ and on the support orbit of $T$.

\begin{definition}[Cumulative squarefree dictionary]
For $d\geq1$, let
\[
  \Dict^G_{n,\leq d}
   =\{q^G_{[T]}: \varnothing\neq T\subseteq\Omega_n,
       \ |T|\leq d\},
\]
with one coordinate for each $G$-orbit of supports.  Its signature map is
\[
  \Fingerprint^G_{n,\leq d}([P]_G)
   =\bigl(q(P)\bigr)_{q\in\Dict^G_{n,\leq d}}.
\]
The reconstruction degree $\rho(n,G)$ is the least $d$ for which this map is
injective.
\end{definition}

The exact-degree coordinate profiles are
\[
\begin{array}{c|c}
G & (|\Dict^G_{n,1}|,|\Dict^G_{n,2}|,\ldots)\\ \hline
S_4&(1,5,13)\\
A_4&(1,7,21)\\
S_5&(1,5,16,61)\\
A_5&(1,6,21,96).
\end{array}
\]
The sums of these rows give the dictionary sizes in the central theorem.  In
particular, a profile entry is an exact-degree count, not a cumulative count.

\subsection{Canonical finite encoding}

The atlas orders arcs lexicographically by $(i,j)$ with $i\neq j$.  A pattern
or support is serialized by the corresponding unsigned bit mask.  Its
canonical representative is the minimum unsigned mask in its $G$-orbit.
Coordinates are ordered by degree and subsequently by canonical support mask.
The identifier therefore contains the action, vertex count, degree, and
canonical mask.  For four vertices, compact table identifiers are resolved
through the coordinate registry; the five-vertex registry uses the mask
directly.  Both forms resolve to a unique canonical support.

This encoding is convenient but not mathematical content: replacing it by
any bijective serialization leaves all orbit sums, signatures, and separator
masks unchanged.  Canonicalization merely permits byte-level comparison of
two independent reconstructions.

\subsection{Orbit incidence and the normalization lock}

For a pattern orbit $R$ and support orbit $C$, let
\[
 A_{R,C}=\#\{T\in C:T\subseteq P\}\quad(P\in R)
\]
and
\[
 M_{R,C}=\#\{P\in R:T\subseteq P\}\quad(T\in C).
\]
Both quantities are well defined by transitivity.  Double-counting incident
pairs $(P,T)\in R\times C$ gives
\begin{equation}
  |R|A_{R,C}=|C|M_{R,C}.
  \label{eq:normalization-lock}
\end{equation}
The signature entries in this paper are the $A_{R,C}$ values of
Equation~\eqref{eq:contained-support-count}.  Equation
\eqref{eq:normalization-lock} is essential: the two incidence normalizations
are not identified unless their orbit-size factors are carried explicitly.
The matrix vocabulary is related to classical orbit-incidence constructions
such as the Kramer--Mesner matrix \cite{KramerMesner1976}, but none of the
finite values below is imported from that literature.

Orbit enumeration and orbit sums sit in a broad classical framework
\cite{Polya1937,Mnukhin1992,RadcliffeScott2006}.  Here the carrier, group,
squarefree support restriction, and integer normalization are fixed before
any separation or optimization is performed.
